\section*{Chapitre \ref{ch:b1:faults-folds-fabrics} --- Failles, plis et fabriques}

\begin{solution}{exo:b1:faults-folds-fabrics:1}
(a) Une \omterm{def:b1:faults-folds-fabrics:fault-types}{faille normale}. (b) Une \omterm{def:b1:faults-folds-fabrics:fault-types}{faille inverse}, presque à pur glissement en \omterm{def:b1:geological-maps:attitude}{pendage}. (c) Une \omterm{def:b1:faults-folds-fabrics:fault-types}{faille
décrochante} \omterm{def:b1:faults-folds-fabrics:fault-types}{senestre}. (d) Une \omterm{def:b1:faults-folds-fabrics:fault-types}{faille décrochante} \omterm{def:b1:faults-folds-fabrics:fault-types}{dextre}. (e) Une \omterm{def:b1:faults-folds-fabrics:fault-types}{faille oblique}~:
le \omterm{def:b1:faults-folds-fabrics:fault-elements}{compartiment supérieur} s'est abaissé et déplacé le long de la direction, vers la gauche~; sa composante
en direction, $\cos\ang{40} = 0.77$ du \omterm{def:b1:faults-folds-fabrics:fault-elements}{glissement net}, est plus grande que sa composante
en \omterm{def:b1:geological-maps:attitude}{pendage}, $0.64$, donc on l'appelle une faille \omterm{def:b1:faults-folds-fabrics:fault-types}{senestre} à composante normale.
(f) Avec un \omterm{def:b1:geological-maps:attitude}{pendage} de \ang{25}, inférieur à \ang{45}, la \omterm{def:b1:faults-folds-fabrics:fault-types}{faille inverse} est un
\omterm{def:b1:faults-folds-fabrics:fault-types}{chevauchement}.
\end{solution}

\begin{solution}{exo:b1:faults-folds-fabrics:2}
Avec $|\sin r| = 1$, le \omterm{def:b1:cross-sections:throw}{rejet vertical} est $300\sin\ang{55} = \qty{245.7}{m}$ et le
\omterm{def:b1:cross-sections:throw}{rejet horizontal} $300\cos\ang{55} = \qty{172.1}{m}$.
\end{solution}

\begin{solution}{exo:b1:faults-folds-fabrics:3}
Pour un \omterm{def:b1:faults-folds-fabrics:fold-classes}{pli droit} dont les \omterm{def:b1:faults-folds-fabrics:fold-elements}{flancs} plongent de $\delta_1$ et $\delta_2$ en sens
opposés, $\iota = \ang{180} - \delta_1 - \delta_2$. (a) \ang{110}~: ouvert.
(b) \ang{35}~: serré. (c) \ang{7}~: isoclinal.
\end{solution}

\begin{solution}{exo:b1:faults-folds-fabrics:4}
Les \omterm{def:b1:faults-folds-fabrics:fault-types}{failles normales} plongent de $\ang{45} + \ang{35}/2 = \ang{62.5}$, les \omterm{def:b1:faults-folds-fabrics:fault-types}{chevauchements} de
$\ang{45} - \ang{17.5} = \ang{27.5}$. Des \omterm{def:b1:faults-folds-fabrics:fault-types}{failles décrochantes} conjuguées font chacune
\ang{27.5} avec $\sigma_1$, de part et d'autre d'elle, donc elles se rencontrent à \ang{55}.
\end{solution}

\begin{solution}{exo:b1:faults-folds-fabrics:5}
(a) La composante en direction est $400\cos\ang{30} = \qty{346.4}{m}$ le long de la
direction, la composante en \omterm{def:b1:geological-maps:attitude}{pendage} $400\sin\ang{-30} = \qty{-200}{m}$, c'est-à-dire
\qty{200}{m} vers le bas du \omterm{def:b1:geological-maps:attitude}{pendage}~; le \omterm{def:b1:cross-sections:throw}{rejet vertical} est $200\sin\ang{50} = \qty{153.2}{m}$ et
le \omterm{def:b1:cross-sections:throw}{rejet horizontal} $200\cos\ang{50} = \qty{128.6}{m}$. (b) La direction est $000$, donc
$\mathbf s = (0, 1, 0)$~; $\mathbf u = (-\cos\ang{50}, 0, \sin\ang{50}) =
(-0.643, 0, 0.766)$. Alors
$\overrightarrow{AA'} = 400(0.866\,\mathbf s - 0.5\,\mathbf u) =
(128.6, 346.4, -153.2)$ en mètres. Sa direction est $\arctan(128.6/346.4) =
\ang{20.4}$ et son \omterm{def:b1:field-lab-methods-1:trend-plunge}{plongement} $\arcsin(153.2/400) = \ang{22.5}$, qui est aussi
$\arcsin(\sin\ang{30}\sin\ang{50})$. (c) Une \omterm{def:b1:faults-folds-fabrics:fault-types}{faille oblique}~: le \omterm{def:b1:faults-folds-fabrics:fault-elements}{pitch} est
entre $0$ et \ang{-90}, plus proche de $0$~: une faille \omterm{def:b1:faults-folds-fabrics:fault-types}{senestre} à composante
normale.
\end{solution}

\begin{solution}{exo:b1:faults-folds-fabrics:6}
(a) Les deux failles font \ang{60} entre elles~; $\sigma_1$ bissecte l'angle
aigu, donc elle est horizontale selon $050$ (et $230$)~; $\sigma_3$ est horizontale
selon $140$~; $\sigma_2$ est verticale, selon la droite où les deux failles verticales
se rencontrent. Chaque faille fait \ang{30} avec $\sigma_1$, donc $\ang{45} - \phi/2 =
\ang{30}$ et $\phi = \ang{30}$. (b) Les deux failles découpent le sol en quatre
coins~; les coins qui contiennent $\sigma_1$, vers $050$ et vers $230$, se déplacent
l'un vers l'autre, et ceux qui contiennent $\sigma_3$ s'écartent, vers $140$
et $320$. Prenons la faille de direction $020$, entre le coin vers $050$ (à son
est) et le coin vers $320$ (à son ouest)~: par rapport au coin ouest,
le coin est se déplace vers $230$ et à l'opposé de $320$, qui ont tous deux une
composante vers $200$ le long de la faille. Un observateur sur le coin ouest,
face à l'est, voit le coin est se déplacer vers $200$, vers la droite~: la faille
de direction $020$ est \omterm{def:b1:faults-folds-fabrics:fault-types}{dextre}. Le même argument, ou la symétrie autour de
$\sigma_1$, rend la faille de direction $080$ \omterm{def:b1:faults-folds-fabrics:fault-types}{senestre}.
\end{solution}

\begin{solution}{exo:b1:faults-folds-fabrics:7}
Les normales vers le haut sont
\begin{align*}
  \mathbf n_1 &= (\sin\ang{40}\sin\ang{120}, \sin\ang{40}\cos\ang{120}, \cos\ang{40}) = (0.557, -0.321, 0.766),\\
  \mathbf n_2 &= (\sin\ang{30}\sin\ang{260}, \sin\ang{30}\cos\ang{260}, \cos\ang{30}) = (-0.492, -0.087, 0.866).
\end{align*}
Leur produit vectoriel est
$(-0.211, -0.859, -0.207)$, qui pointe vers le bas~: la \omterm{def:b1:faults-folds-fabrics:fold-elements}{ligne de charnière} plonge de
$\arcsin 0.207/0.909 = \ang{13.1}$ vers $\ang{180} + \arctan(0.211/0.859) =
\ang{193.8}$. $\mathbf n_1\cdot\mathbf n_2 = 0.417$, l'angle entre les
normales est \ang{65.3}, et $\iota = \ang{114.7}$~: un \omterm{def:b1:faults-folds-fabrics:fold-classes}{pli ouvert}.
\end{solution}

\begin{solution}{exo:b1:faults-folds-fabrics:8}
$\mathbf n_1 = (0.321, 0.383, 0.866)$ pour la stratification et
$\mathbf n_2 = (0.337, -0.925, 0.174)$ pour le \omterm{def:b1:faults-folds-fabrics:foliation}{clivage}. Leur produit vectoriel,
pris vers le bas, donne une droite qui plonge de \ang{25.4} vers \ang{074.8}. D'après la
\cref{prop:b1:faults-folds-fabrics:intersection}, c'est la direction de la
\omterm{def:b1:faults-folds-fabrics:fold-elements}{ligne de charnière}~: le \omterm{def:b1:geological-maps:fold}{pli} plonge d'environ \ang{25} vers l'est-nord-est, ce que
tout autre \omterm{def:b1:geological-maps:geological-map}{affleurement} de l'un ou l'autre flanc devrait confirmer.
\end{solution}

\begin{solution}{exo:b1:faults-folds-fabrics:9}
(a) Un observateur au sud de la faille, face au nord, voit la partie nord du
dyke déplacée vers l'est, vers la droite~: la faille est \omterm{def:b1:faults-folds-fabrics:fault-types}{dextre}. (b) Le dyke est
un plan vertical perpendiculaire à la faille, donc son décalage selon la direction est
la composante en direction~: \qty{350}{m}, qui est tout le glissement si les
\omterm{def:b1:faults-folds-fabrics:fault-rocks}{stries} sont horizontales. (c) Le glissement fait alors \ang{20} avec la direction
dans le plan de faille vertical~: sa composante horizontale est toujours \qty{350}{m},
donc le \omterm{def:b1:faults-folds-fabrics:fault-elements}{glissement net} est $350/\cos\ang{20} = \qty{372.5}{m}$ et sa composante
verticale $350\tan\ang{20} = \qty{127.4}{m}$. Le bloc nord s'est déplacé vers l'est,
et les \omterm{def:b1:faults-folds-fabrics:fault-rocks}{stries} descendent vers l'est, donc il est aussi descendu, de
\qty{127.4}{m}. Un dyke vertical ne peut montrer ce mouvement vertical, puisque
le faire glisser le long de lui-même laisse sa trace inchangée.
\end{solution}

\begin{solution}{exo:b1:faults-folds-fabrics:10}
(a) La normale à la faille fait $\ang{90} - \theta$ avec $\sigma_1$, donc d'après la
\cref{prop:b1:stress-strain-rheology:mohr},
$\sigma_n' = \sigma_1'\sin^2\theta + \sigma_3'\cos^2\theta$ et
$\tau = (\sigma_1' - \sigma_3')\sin\theta\cos\theta$. Le glissement commence quand
$\tau = \mu\sigma_n'$~; en divisant par $\sigma_3'\cos^2\theta$ et en posant
$R = \sigma_1'/\sigma_3'$, $(R - 1)\tan\theta = \mu(R\tan^2\theta + 1)$, donc
$R\tan\theta(1 - \mu\tan\theta) = \tan\theta + \mu$, qui est la formule.
(b) Avec $\sigma_1$ verticale, $\theta$ est \ang{90} moins le \omterm{def:b1:geological-maps:attitude}{pendage}~:
$\theta = \ang{30}$ donne $R = 3.12$~; $\theta = \ang{45}$, $R = 4.00$~;
$\theta = \ang{55}$, $R = 9.92$. (c) Le dénominateur s'annule quand
$\tan\theta = 1/\mu$, $\theta = \ang{59.0}$~; au-delà, $R$ serait négatif~:
aucun rapport de \omterm{def:b1:stress-strain-rheology:lithostatic}{contraintes effectives} positives ne peut faire glisser la faille, si grande que devienne
$\sigma_1'$. Les \omterm{def:b1:faults-folds-fabrics:fault-types}{failles normales} de \omterm{def:b1:geological-maps:attitude}{pendage} inférieur à \ang{31.0} ne peuvent glisser
sous une $\sigma_1$ verticale avec $\mu = 0.6$. (d) Le plan le mieux orienté est à
$\ang{45} - \phi/2$ de $\sigma_1$. Pour $\mu = 0.6$, $\phi = \ang{31.0}$~:
les \omterm{def:b1:faults-folds-fabrics:fault-types}{failles normales} plongent de \ang{60.5}, les \omterm{def:b1:faults-folds-fabrics:fault-types}{chevauchements} de \ang{29.5}. Pour $\mu = 0.85$,
$\phi = \ang{40.4}$~: \ang{65.2} et \ang{24.8}.
\end{solution}

\begin{solution}{exo:b1:faults-folds-fabrics:11}
En regardant vers le nord, l'est est à droite, et les plans plongeant vers l'ouest montent vers la
droite. (a) En B la stratification plonge dans le même sens que le \omterm{def:b1:faults-folds-fabrics:foliation}{clivage}, plus
fortement~: d'après la \cref{prop:b1:faults-folds-fabrics:cleavage-bedding}(a), B est sur un
flanc renversé~; A et C, où la stratification est moins raide, sont à
l'endroit. (b) Le \omterm{def:b1:faults-folds-fabrics:foliation}{clivage} monte à \ang{60}~; en A et C la stratification monte à
\ang{40}, et la rotation aiguë de l'un à l'autre est horaire, de
\ang{20}, déplaçant le haut du \omterm{def:b1:faults-folds-fabrics:foliation}{clivage} vers l'est~; en B la stratification monte
à \ang{80} et la rotation est antihoraire, de \ang{20}, le déplaçant vers
l'ouest. (c) Le sens change entre A et B et de nouveau entre B et C~:
il y a une \omterm{def:b1:faults-folds-fabrics:fold-elements}{charnière} dans chaque intervalle. En A, qui est à l'endroit, le haut du
\omterm{def:b1:faults-folds-fabrics:foliation}{clivage} se déplace vers l'est, vers l'\omterm{def:b1:geological-maps:fold}{anticlinal} au sommet du flanc~: la \omterm{def:b1:faults-folds-fabrics:fold-elements}{charnière}
entre A et B est un \omterm{def:b1:geological-maps:fold}{anticlinal}, et celle entre B et C un \omterm{def:b1:geological-maps:fold}{synclinal}.
(d) Dans la notation de la démonstration, les \omterm{def:b1:faults-folds-fabrics:fold-elements}{flancs} à l'endroit plongent de
$\gamma - \iota/2 = \ang{40}$ et les \omterm{def:b1:faults-folds-fabrics:fold-elements}{flancs} renversés de $\gamma + \iota/2 =
\ang{80}$, donc $\gamma = \ang{60}$, comme le dit le \omterm{def:b1:faults-folds-fabrics:foliation}{clivage}, et $\iota =
\ang{40}$~: des \omterm{def:b1:faults-folds-fabrics:fold-classes}{plis serrés}. Les plans axiaux sont de direction nord et plongent de \ang{60} vers
l'ouest, \omterm{def:b1:geological-maps:attitude}{attitude} $270/60$. Les \omterm{def:b1:faults-folds-fabrics:fold-elements}{flancs} courts ont dépassé la verticale avec
leurs sommets vers l'est~: les \omterm{def:b1:geological-maps:fold}{plis} ont une \omterm{def:b1:faults-folds-fabrics:vergence}{vergence} vers l'est, comme sur la figure
du chapitre.
\end{solution}

\begin{solution}{exo:b1:faults-folds-fabrics:12}
(a) Dans un \omterm{def:b1:geological-maps:fold}{pli} en chevron symétrique chaque flanc fait l'angle $\iota/2$ avec le
\omterm{def:b1:faults-folds-fabrics:fold-elements}{plan axial}, qui est perpendiculaire à la direction du raccourcissement~; un flanc de
longueur $l$ couvre donc $l\sin(\iota/2)$ à travers les \omterm{def:b1:geological-maps:fold}{plis}, et toute la
couche $L_0\sin(\iota/2)$. (b) $\sin\ang{30} = 0.5$~: un raccourcissement de
\qty{50}{\percent}. (c) $\sin(\iota/2) = 0.7$ donne $\iota = \ang{88.9}$.
(d) $\iota = \ang{156}$ et $\sin\ang{78} = 0.978$~: un raccourcissement de
\qty{2.2}{\percent}, qui est aussi $1 - \cos\ang{12}$.
\end{solution}

\begin{solution}{pb:b1:faults-folds-fabrics:1}
\textbf{1.} $\mathbf n = (\sin\delta\sin\alpha, \sin\delta\cos\alpha,
\cos\delta)$ donne
\[
  \mathbf n_1 = (0.20855, 0.03415, 0.97742), \qquad
  \mathbf n_2 = (-0.20855, 0.03415, 0.97742).
\]

\textbf{2.} $\mathbf n_1\cdot\mathbf n_2 = -0.04349 + 0.00117 + 0.95535 =
0.91302$, et l'angle entre les normales est \ang{24.1}.

\textbf{3.} D'après la \cref{prop:b1:faults-folds-fabrics:hinge},
$\iota = \ang{180} - \ang{24.1} = \ang{155.9}$.

\textbf{4.} Entre \ang{170} et \ang{90}~: un \omterm{def:b1:faults-folds-fabrics:fold-classes}{pli ouvert}.

\textbf{5.} Les \omterm{def:b1:faults-folds-fabrics:fold-elements}{flancs} plongent vers $080.7$ et $279.3$, non plein est et ouest,
parce que le \omterm{def:b1:geological-maps:fold}{pli} plonge~: leurs lignes de plus grande pente ne sont pas dans le profil, qui
ne montre qu'un \omterm{def:b1:cross-sections:apparent-dip}{pendage apparent}. Dans une coupe est--ouest, d'après le
\cref{thm:b1:cross-sections:apparent-dip}, $\tan\beta = \tan\ang{12.2}
\cos\ang{9.3} = 0.2134$, donc $\beta = \ang{12.0}$, et $\ang{180} -
2\times\ang{12.0} = \ang{155.9}$ (le profil est incliné de \ang{2} par rapport à cette
coupe, ce qui ne change rien à cette précision).

\textbf{6.} Ses \omterm{def:b1:faults-folds-fabrics:fold-elements}{flancs} sont parallèles à ceux de l'\omterm{def:b1:geological-maps:fold}{anticlinal}, donc son \omterm{def:b1:faults-folds-fabrics:fold-elements}{angle entre
flancs} est aussi \ang{155.9}.

\textbf{7.} $\mathbf n_1\times\mathbf n_2 = (0.03415\times0.97742 -
0.97742\times0.03415,\; 0.97742\times(-0.20855) - 0.20855\times0.97742,\;
0.20855\times0.03415 + 0.03415\times0.20855) = (0, -0.40768, 0.01424)$.
Le vecteur orienté vers le bas est $(0, 0.40768, -0.01424)$~: la \omterm{def:b1:faults-folds-fabrics:fold-elements}{ligne de charnière} est de direction $000$ et
plonge de $\arctan(0.01424/0.40768) = \ang{2.0}$ vers le nord.

\textbf{8.} Le \omterm{def:b1:faults-folds-fabrics:fold-elements}{plan axial} contient la \omterm{def:b1:faults-folds-fabrics:fold-elements}{ligne de charnière} et bissecte les deux \omterm{def:b1:faults-folds-fabrics:fold-elements}{flancs}~;
sa normale est selon $\mathbf n_1 - \mathbf n_2 = (0.41710, 0, 0)$, horizontale
et est--ouest. Il est vertical et de direction nord--sud~: $090/90$.
L'\omterm{def:b1:geological-maps:fold}{anticlinal} est droit.

\textbf{9.} Un \omterm{def:b1:faults-folds-fabrics:fold-classes}{pli ouvert}, droit, plongeant de \ang{2.0} vers le nord.

\textbf{10.} D'après la \cref{prop:b1:geological-maps:folds}, les \omterm{def:b1:geological-maps:geological-map}{affleurements} d'un
\omterm{def:b1:geological-maps:fold}{anticlinal} se ferment dans la direction du \omterm{def:b1:field-lab-methods-1:trend-plunge}{plongement}~: vers le nord. Son cœur contient les
roches les plus anciennes~: le grès à la base de la couverture et, là où la rivière a
creusé le plus profond, le socle de schiste et de granite en dessous.

\textbf{11.} De \omterm{def:b1:faults-folds-fabrics:fold-elements}{plan axial}, le \omterm{def:b1:faults-folds-fabrics:foliation}{clivage} est vertical et de direction nord--sud,
$090/90$. D'après la \cref{prop:b1:faults-folds-fabrics:intersection}, sa
\omterm{def:b1:faults-folds-fabrics:lineation}{linéation d'intersection} avec la stratification est parallèle à la \omterm{def:b1:faults-folds-fabrics:fold-elements}{ligne de charnière}~: direction
$000$, \omterm{def:b1:field-lab-methods-1:trend-plunge}{plongement} \ang{2.0}. Directement~: le plan $080.7/12.2$ rencontre le plan
vertical de direction nord selon une droite de \omterm{def:b1:field-lab-methods-1:trend-plunge}{plongement} $\arctan(\tan\ang{12.2}
\cos\ang{80.7}) = \ang{2.0}$ vers $000$.

\textbf{12.} En regardant vers le nord, le flanc est monte vers la gauche de \ang{12.0}
et le \omterm{def:b1:faults-folds-fabrics:foliation}{clivage} est vertical~: la rotation aiguë du \omterm{def:b1:faults-folds-fabrics:foliation}{clivage} à la
stratification, de \ang{78.0}, est antihoraire et déplace le haut du
\omterm{def:b1:faults-folds-fabrics:foliation}{clivage} vers l'ouest. Le flanc est à l'endroit, donc d'après la
\cref{prop:b1:faults-folds-fabrics:cleavage-bedding}(b) la \omterm{def:b1:faults-folds-fabrics:fold-elements}{charnière} de l'\omterm{def:b1:geological-maps:fold}{anticlinal} est
à l'ouest, en $x = \qty{3.5}{km}$, comme le montre la carte.

\textbf{13.} La direction est $120 - 90 = 030$~; le plan plonge vers $120$, donc
le \omterm{def:b1:faults-folds-fabrics:fault-elements}{compartiment supérieur} se trouve du côté est-sud-est de la trace.

\textbf{14.} La faille est normale~: $\sigma_1$ était verticale.

\textbf{15.} $\sigma_2$ est horizontale selon la direction, $030$--$210$~;
$\sigma_3$ horizontale en travers, $120$--$300$.

\textbf{16.} $\ang{45} + \phi/2 = \ang{60}$ donne $\phi = \ang{30}$, et
$\mu = \tan\ang{30} = 0.577$.

\textbf{17.} La famille conjuguée est aussi de direction $030$ et plonge de \ang{60} dans l'autre
sens~: $300/60$. Les deux familles se rencontrent à $\ang{180} - 2\times\ang{60} = \ang{60}$,
avec $\sigma_1$ verticale bissectant l'angle aigu.

\textbf{18.} Avec $\mu = 0.85$, $\phi = \ang{40.4}$ et les \omterm{def:b1:faults-folds-fabrics:fault-types}{failles normales} les mieux orientées
plongent de \ang{65.2}~; la faille de Varden, à \ang{60}, en est à \ang{5.2}.
D'après l'\cref{exo:b1:faults-folds-fabrics:10}, elle glisse au rapport
$\sigma_1'/\sigma_3' = (1 + 0.85\cot\ang{30})/(1 - 0.85\tan\ang{30}) = 4.85$,
contre $4.68$ pour le meilleur plan~: elle est presque, mais pas exactement, la mieux
orientée, et glisserait de nouveau à un rapport de contraintes supérieur de \qty{4}{\percent}
seulement.

\textbf{19.} Le glissement est vers le bas du \omterm{def:b1:geological-maps:attitude}{pendage}~: $r = \ang{-90}$~; le \omterm{def:b1:faults-folds-fabrics:fault-elements}{vecteur glissement} est de direction
$120$ et plonge de \ang{60}.

\textbf{20.} $(\cos\ang{60}\sin\ang{120}, \cos\ang{60}\cos\ang{120},
-\sin\ang{60}) = (0.4330, -0.2500, -0.8660)$.

\textbf{21.} $e = 200/\tan\ang{60} = \qty{115.5}{m}$.

\textbf{22.} Est~: $115.47\sin\ang{120} = \qty{100.0}{m}$~; nord~:
$115.47\cos\ang{120} = \qty{-57.7}{m}$, c'est-à-dire \qty{57.7}{m} vers le sud.

\textbf{23.} $L = 200/\sin\ang{60} = \qty{230.9}{m}$, et
$200^2 + 115.47^2 = 40\,000 + 13\,333 = 53\,333 = 230.94^2$. Le \omterm{def:b1:cross-sections:throw}{rejet vertical} et le
\omterm{def:b1:cross-sections:throw}{rejet horizontal} sont les composantes verticale et horizontale du glissement, qui est dans
le plan vertical du \omterm{def:b1:geological-maps:attitude}{pendage}~: ce sont les côtés d'un triangle rectangle dont
l'hypoténuse est le glissement.

\textbf{24.} Au mètre près, le \omterm{def:b1:faults-folds-fabrics:fault-elements}{glissement net} sur la faille de Varden est de
\end{solution}
